%! AP Statistics — Unit 4, Lesson 4.7 — Group Activity
\documentclass[10pt]{article}
\usepackage{apstats-article}
\usepackage{apstats-boxes}

\begin{document}

\pageheader{Unit 4, Lesson 4.7}{Group Activity: Building and Reading Distributions}
\namepartnerperiod

\vspace{0.1in}

\begin{tcolorbox}[colback=sky!60, colframe=navy, title=\textbf{Tier R --- Remediate},
    arc=2mm, boxrule=0.8pt, left=4mm, right=4mm, top=3mm, bottom=3mm]
\small
A survey of 50 students asked how many siblings each student has. Results:
\begin{center}
\renewcommand{\arraystretch}{1.4}
\begin{tabular}{@{}cccc@{}}
\textbf{Siblings ($x$)} & \textbf{Frequency} & \textbf{$P(X=x)$} & \textbf{$P(X \le x)$} \\
\hline
0 & 8  & \blank{2cm} & \blank{2cm} \\
1 & 20 & \blank{2cm} & \blank{2cm} \\
2 & 15 & \blank{2cm} & \blank{2cm} \\
3 & 7  & \blank{2cm} & \blank{2cm} \\
\end{tabular}
\end{center}

\begin{enumerate}[leftmargin=1.5em, itemsep=8pt]
  \item Fill in the $P(X=x)$ column. Divide each frequency by the total.

        \vspace{0.1in}

  \item Do these probabilities form a valid distribution? Check both conditions.\\
        \writeline

  \item Fill in the $P(X \le x)$ column.

        \vspace{0.05in}

  \item Using your table, find $P(X \ge 2)$. Show your work.

        \vspace{0.3in}

  \item Is it more likely that a randomly chosen student has 0 siblings or 3 siblings?
        Explain.\\
        \writeline
\end{enumerate}
\end{tcolorbox}

\vspace{0.12in}

\begin{tcolorbox}[colback=goldbg!60, colframe=goldacc, title=\textbf{Tier A --- Approaching Proficiency},
    arc=2mm, boxrule=0.8pt, left=4mm, right=4mm, top=3mm, bottom=3mm]
\small
$X$ = number of books a student reads during summer break. The probability distribution is:
\begin{center}
\renewcommand{\arraystretch}{1.4}
\begin{tabular}{@{}cccccc@{}}
$x$ & 0 & 1 & 2 & 3 & 4 \\
\hline
$P(X=x)$ & 0.05 & 0.20 & 0.35 & 0.25 & 0.15 \\
\end{tabular}
\end{center}

\begin{enumerate}[leftmargin=1.5em, itemsep=8pt]
  \item Verify this is a valid probability distribution.\\
        \writeline

  \item Find $P(1 \le X \le 3)$. Interpret your answer in context.\\
        \writeline\\[1pt]\writeline

  \item Find $P(X \ge 3)$. Interpret your answer in context.\\
        \writeline\\[1pt]\writeline

  \item Describe the shape of this distribution. Where is it centered?\\
        \writeline\\[1pt]\writeline

  \item Write an interpretation of this distribution in context (shape, center, spread).\\
        \writeline\\[1pt]\writeline
\end{enumerate}
\end{tcolorbox}

\vspace{0.12in}

\begin{tcolorbox}[colback=greenbg!60, colframe=greenacc, title=\textbf{Tier E --- Extension},
    arc=2mm, boxrule=0.8pt, left=4mm, right=4mm, top=3mm, bottom=3mm]
\small
Two classes were surveyed about the number of after-school activities they participate in.

\textbf{Class A:}
\begin{center}
\renewcommand{\arraystretch}{1.3}
\begin{tabular}{@{}ccccc@{}}
$x$ & 0 & 1 & 2 & 3 \\
\hline
$P(X=x)$ & 0.40 & 0.35 & 0.20 & 0.05 \\
\end{tabular}
\end{center}

\textbf{Class B:}
\begin{center}
\renewcommand{\arraystretch}{1.3}
\begin{tabular}{@{}ccccc@{}}
$x$ & 0 & 1 & 2 & 3 \\
\hline
$P(X=x)$ & 0.10 & 0.25 & 0.40 & 0.25 \\
\end{tabular}
\end{center}

\begin{enumerate}[leftmargin=1.5em, itemsep=8pt]
  \item Verify both distributions are valid.

        \vspace{0.2in}

  \item Compare the two distributions: which class tends to participate in \emph{more}
        activities? Support your answer using probabilities.\\
        \writeline\\[1pt]\writeline

  \item For each class, find $P(X \ge 2)$ and compare.

        \vspace{0.25in}

  \item If a student is chosen at random from Class A, what is the probability they
        participate in \emph{at least} as many activities as the most common value in
        Class B? Show your work.\\
        \writeline\\[1pt]\writeline

  \item Describe the difference in shape between the two distributions.\\
        \writeline\\[1pt]\writeline
\end{enumerate}
\end{tcolorbox}

\end{document}
